%% Question bank chapter file - generated by bank_tools.py build (do not edit while the bank is still being built:
%% it is regenerated from the per-test data in Question Bank\data). Every question is one bankq block:
%%   \begin{bankq}[id=..., source={...}, paper=.., qno=.., type=..., subtopic={...}, answer={...}, answer_src=key|solved,
%%                 working={...}, difficulty=1-10, diff_parts={K S I C F}, diff_note={...},
%%                 group=<passage id or empty>, also_in={duplicate copies in other tests}]
%% difficulty = min(10, K+S+I+C+F-1): K concepts 1-3, S steps 1-4, I insight 0-2, C computation 0-2, F format 0-1
%%   ...question LaTeX, options, figures (\bfig / \bfigi -> img/)...
%%   \end{bankq}
%% Passages shared by several questions are bankpassage blocks; their questions carry group=<passage id>.
\documentclass[10pt]{article}
\usepackage[a4paper,margin=1.8cm]{geometry}
\usepackage[T1]{fontenc}
\usepackage{amsmath,amssymb,mathtools}
\usepackage{graphicx}
\usepackage[export]{adjustbox}
\usepackage{xcolor,needspace,keyval}
\usepackage{enumitem}
\usepackage{booktabs,array,tabularx}
\usepackage[hidelinks]{hyperref}
\setlength{\parindent}{0pt}
\setlength{\parskip}{3pt}
\makeatletter
\define@key{bq}{id}{\def\bq@id{#1}}\define@key{bq}{source}{\def\bq@src{#1}}\define@key{bq}{paper}{\def\bq@paper{#1}}
\define@key{bq}{qno}{\def\bq@qno{#1}}\define@key{bq}{type}{\def\bq@type{#1}}\define@key{bq}{subtopic}{\def\bq@sub{#1}}
\define@key{bq}{answer}{\def\bq@ans{#1}}\define@key{bq}{answer_src}{\def\bq@asrc{#1}}\define@key{bq}{working}{\def\bq@work{#1}}
\define@key{bq}{group}{\def\bq@group{#1}}\define@key{bq}{also_in}{\def\bq@also{#1}}
\define@key{bq}{difficulty}{\def\bq@diff{#1}}\define@key{bq}{diff_parts}{\def\bq@dparts{#1}}\define@key{bq}{diff_note}{\def\bq@dnote{#1}}
\newenvironment{bankq}[1][]{\def\bq@work{}\def\bq@also{}\def\bq@group{}\def\bq@diff{?}\def\bq@dparts{}\def\bq@dnote{}\setkeys{bq}{#1}%
  \par\needspace{4\baselineskip}\medskip\noindent{\small\textbf{\bq@id}\hfill\color{gray}[\bq@src, Paper \bq@paper, Q\bq@qno]
  \bq@type\ $\cdot$ \bq@sub\ $\cdot$ \textbf{D\bq@diff}}\par\nopagebreak\relax}
  {\par{\small\color{teal}Answer: \textbf{\bq@ans}\ifx\bq@work\@empty\else\ (solved: \bq@work)\fi
  \ifx\bq@also\@empty\else\par\color{gray}Also in: \bq@also\fi}\par}
\define@key{bp}{id}{\def\bp@id{#1}}\define@key{bp}{source}{\def\bp@src{#1}}\define@key{bp}{questions}{\def\bp@qs{#1}}
\newenvironment{bankpassage}[1][]{\setkeys{bp}{#1}\par\needspace{4\baselineskip}\bigskip\noindent
  {\small\textbf{Passage \bp@id}\hfill\color{gray}[\bp@src; questions \bp@qs]}\par\nopagebreak\relax\itshape}{\par}
\makeatother
\newcommand{\bfig}[2]{\par\smallskip{\centering\includegraphics[width=#1pt,max width=0.62\linewidth,max height=0.32\textheight]{img/#2}\par}\smallskip}
\newcommand{\bfigi}[2]{\includegraphics[width=#1pt,max width=\linewidth,max height=0.32\textheight,valign=c]{img/#2}}
\newcommand{\bmissing}[1]{\textit{[figure missing in the source conversion: \detokenize{#1}]}}
\newenvironment{qoptions}{\begin{itemize}[leftmargin=1.8em,itemsep=1pt,topsep=2pt,label={}]}{\end{itemize}}
\newcommand{\opt}[2]{\item[(#1)] #2}
\providecommand{\fwstart}{}\providecommand{\fwend}{}\providecommand{\qfullsplit}{}
\newcolumntype{L}[1]{>{\raggedright\arraybackslash}p{#1}}
\newcolumntype{C}[1]{>{\centering\arraybackslash}p{#1}}
\newcolumntype{R}[1]{>{\raggedleft\arraybackslash}p{#1}}
\providecommand{\sectionhdr}[1]{}\providecommand{\subsectionhdr}[1]{}


\begin{document}
{\Large\bfseries Gaseous State --- Question Bank}\par\smallskip
{\small Chemistry chapter C03. Questions from all converted test papers, grouped by source series. Answers: printed key where it matches the paper, otherwise solved (answer\_src=solved, with working). Exact repeats across papers appear once, with also\_in listing the other copies.}\par\medskip\hrule

\section*{TOAS papers}
\begin{bankq}[id=TOAS04-P1-CHEM-S2-Q07, source={TOAS - Test 04 (15 Mar 2022)}, paper=1, qno=7, type=multi, subtopic={Gas laws}, answer={A,B,C}, answer_src=key, difficulty=5, diff_parts={K1 S3 I0 C1 F1}, diff_note={concepts: gas mixtures and pistons | steps: moles of each gas; equal final pressures; position of piston A; position of piston B; pressures}]
Equal mass of three gases He, $\mathrm{O_2}$ \& $\mathrm{SO_2}$ are present in a container initially as shown -
\bfig{166.5}{TOAS04_p1_chem_s2_q07_1.png}
If mechanical stoppers are removed and movable pistons are allowed to attain equilibrium -

Choose the \textbf{incorrect} statement -
\begin{qoptions}
\opt{A}{Position of piston $A$ from the left end of the container is 15 cm}
\opt{B}{Position of piston $B$ from the left end of the container is 7 cm}
\opt{C}{Final pressure of $\mathrm{O_2}$ gas is 10 atm}
\opt{D}{Final pressure of He gas is 20 atm}
\end{qoptions}
\end{bankq}
\begin{bankq}[id=TOAS06-P1-CHEM-S1-Q02, source={TOAS - Test 06 (17 Apr 2022)}, paper=1, qno=2, type=single, subtopic={Real gases}, answer={A}, answer_src=key, difficulty=3, diff_parts={K1 S2 I1 C0 F0}, diff_note={concepts: virial equation | steps: substitute 1/V\_m in terms of p; match coefficients of p\^{}2; express C' | I: invert the series}]
Virial equation of state can be written in terms of molar volume ($V_m$) or in terms of pressure as
\[ Z=1+\frac{B}{V_m}+\frac{C}{V_m^2}+\frac{D}{V_m^3}+\ldots\ldots\ (1) \]
or \quad $Z=1+B'p+C'P^2+\ldots\ldots$ (2)\\
Find $C'$ in terms of $B$ and $C$ at temperature $T$.
\begin{qoptions}
\opt{A}{$C'=\dfrac{C-B^2}{(RT)^2}$}
\opt{B}{$C'=\dfrac{C-B}{(RT)^2}$}
\opt{C}{$C'=\dfrac{C}{RT}$}
\opt{D}{$C'=\dfrac{C}{(RT)^2}$}
\end{qoptions}
\end{bankq}
\begin{bankq}[id=TOAS09-P1-CHEM-S2-Q09, source={TOAS - Test 09 (30 May 2022)}, paper=1, qno=9, type=multi, subtopic={Gas laws}, answer={A}, answer_src=key, difficulty=3, diff_parts={K1 S1 I1 C0 F1}, diff_note={concepts: distribution of molecular speeds | steps: effect of temperature and molar mass on the curves; statements (i) to (iv) | I: area under each distribution is fixed}]
\bfig{183.4}{TOAS09_p1_chem_s2_q09_1.png}
Which of the following statement is correct.\par
(i) If A \& B are same gases then T$_2>$ T$_1$\par
(ii) If T$_1$ = T$_2$ then A$_{(g)}$ may be SO$_{2(g)}$ \& B$_{(g)}$ may be CH$_{4(g)}$\par
(iii) If B \& C are same gases then T$_3>$ T$_2$\par
(iv) If T$_1$ = T$_2$ = T$_3$ then A$_{(g)}$ may be CH$_4$, B$_{(g)}$ may be SO$_2$ \& C$_{(g)}$ may be SO$_{3(g)}$
\begin{qoptions}
\opt{A}{i, ii \& iii only}
\opt{B}{i, ii, \& iv only}
\opt{C}{All are correct}
\opt{D}{ii \& iii only}
\end{qoptions}
\end{bankq}
\begin{bankq}[id=TOAS09-P2-CHEM-S1-Q18, source={TOAS - Test 09 (30 May 2022)}, paper=2, qno=18, type=single, subtopic={Kinetic theory}, answer={D}, answer_src=key, difficulty=2, diff_parts={K1 S1 I1 C0 F0}, diff_note={concepts: effusion of a gas mixture | steps: Graham's law and the mole fraction of the escaping gas; solve P | I: composition of effused gas}]
A vessel contains He gas at 20 atm and CH$_4$ gas at `P' atm. When a pin-hole is made in the vessel, the gases coming out initially contains 80\% He gas, by mole. The value of `P' is :-
\begin{qoptions}
\opt{A}{1 atm}
\opt{B}{40 atm}
\opt{C}{2.5 atm}
\opt{D}{10 atm}
\end{qoptions}
\end{bankq}
\begin{bankq}[id=TOAS09-P2-CHEM-S1-Q19, source={TOAS - Test 09 (30 May 2022)}, paper=2, qno=19, type=single, subtopic={Real gases}, answer={A}, answer_src=key, difficulty=1, diff_parts={K1 S1 I0 C0 F0}, diff_note={concepts: critical and Boyle temperatures | steps: T\_B = 27 T\_c/8 for a van der Waals gas; Tc = 120 K}]
A real gas behaves like an ideal gas at 132$^{\circ}$C in low pressure range. Critical temperature for this gas is -
\begin{qoptions}
\opt{A}{120 K}
\opt{B}{39 K}
\opt{C}{132 K}
\opt{D}{15 K}
\end{qoptions}
\end{bankq}
\begin{bankq}[id=TOAS12-P2-CHEM-S1-Q03, source={TOAS - Test 12 (10 Jul 2022)}, paper=2, qno=3, type=single, subtopic={Real gases}, answer={D}, answer_src=key, difficulty=3, diff_parts={K1 S2 I0 C1 F0}, diff_note={concepts: van der Waals constant from compressibility factor | steps: Z = PV/RT gives the molar volume; use the van der Waals equation; solve a}]
The compressibility factor of a Vanderwaal gas is 0.5 at 27$^\circ$C and 24 atm. The value of Vanderwaal constant `a' (in atm-L$^2$/mol$^2$) is\par
($b=0.10$ L/mol, $R=0.08$ L-atm/K-mol)
\begin{qoptions}
\opt{A}{2.4}
\opt{B}{4.0}
\opt{C}{3.6}
\opt{D}{9.0}
\end{qoptions}
\end{bankq}
\begin{bankq}[id=TOAS16-P2-CHEM-S1-Q04, source={TOAS - Test 16 (04 Sep 2022)}, paper=2, qno=4, type=multi, subtopic={Real gases}, answer={B,C,D}, answer_src=key, difficulty=4, diff_parts={K1 S2 I1 C0 F1}, diff_note={concepts: compressibility of real gases | steps: repulsion and attraction dominance; effect of temperature relative to Boyle point; pressure range}]
Compressibility of real gas will be less than ideal gas :
\begin{qoptions}
\opt{A}{At low pressure when, temperature of gas $<$ Boyle's temperature}
\opt{B}{At low pressure when, temperature of gas $>$ Boyle's temperature}
\opt{C}{At very high pressure when, temperature of gas $<$ Boyle's temperature}
\opt{D}{At very high pressure when, temperature of gas $>$ Boyle's temperature}
\end{qoptions}
\end{bankq}
\begin{bankq}[id=TOAS19-P1-CHEM-S2-Q09, source={TOAS - Test 19 (05 Feb 2023)}, paper=1, qno=9, type=single, subtopic={Real gases}, answer={D}, answer_src=key, difficulty=1, diff_parts={K1 S1 I0 C0 F0}, diff_note={concepts: compressibility factor of a van der Waals gas | steps: expressions in different regimes; match the descriptions}]
\begin{center}\footnotesize\setlength{\tabcolsep}{2pt}
\begin{tabular}{|L{0.8cm}L{2.2cm}|L{0.6cm}L{2.2cm}|}\hline
\multicolumn{2}{|L{3.0cm}|}{\centering List-I\\(condition of vander walls gas)}&\multicolumn{2}{L{2.8cm}|}{\centering List-II\\(compressibility factor)}\\\hline
(I)&At low to moderate pressure&(P)&$Z=1+\dfrac{Pb}{RT}$\\\hline
(II)&At critical condition of $P,V,T$&(Q)&$Z=1-\dfrac{a}{V_mRT}$\\\hline
(III)&For H$_2$ \& helium gas at room temperature&(R)&$Z=3/8$\\\hline
(IV)&At very low pressure \& high temperature&(S)&$Z=1$\\\hline
&&(T)&Positive derivation from ideality\\\hline
&&(U)&Negative derivation from ideality\\\hline
\end{tabular}\end{center}
Select the CORRECT option -
\begin{qoptions}
\opt{A}{III $\to$ Q}
\opt{B}{II $\to$ S}
\opt{C}{IV $\to$ P}
\opt{D}{III $\to$ P}
\end{qoptions}
\end{bankq}
\begin{bankq}[id=TOAS19-P1-CHEM-S2-Q10, source={TOAS - Test 19 (05 Feb 2023)}, paper=1, qno=10, type=single, subtopic={Real gases}, answer={B}, answer_src=key, difficulty=1, diff_parts={K1 S1 I0 C0 F0}, diff_note={concepts: compressibility factor of a van der Waals gas | steps: expressions in different regimes; match the descriptions}]
\begin{center}\footnotesize\setlength{\tabcolsep}{2pt}
\begin{tabular}{|L{0.8cm}L{2.2cm}|L{0.6cm}L{2.2cm}|}\hline
\multicolumn{2}{|L{3.0cm}|}{\centering List-I\\(condition of vander walls gas)}&\multicolumn{2}{L{2.8cm}|}{\centering List-II\\(compressibility factor)}\\\hline
(I)&At low to moderate pressure&(P)&$Z=1+\dfrac{Pb}{RT}$\\\hline
(II)&At critical condition of $P,V,T$&(Q)&$Z=1-\dfrac{a}{V_mRT}$\\\hline
(III)&For H$_2$ \& helium gas at room temperature&(R)&$Z=3/8$\\\hline
(IV)&At very low pressure \& high temperature&(S)&$Z=1$\\\hline
&&(T)&Positive derivation from ideality\\\hline
&&(U)&Negative derivation from ideality\\\hline
\end{tabular}\end{center}
Select the CORRECT option -
\begin{qoptions}
\opt{A}{I $\to$ P}
\opt{B}{II $\to$ R}
\opt{C}{IV $\to$ T}
\opt{D}{IV $\to$ U}
\end{qoptions}
\end{bankq}

\section*{TOPS papers}
\begin{bankq}[id=TOPS02-P2-CHEM-S3-Q03, source={TOPS - Test 02 (28 Aug 2022)}, paper=2, qno=3, type=numerical, subtopic={Gas laws}, answer={1.00}, answer_src=key, difficulty=6, diff_parts={K2 S2 I1 C1 F1}, diff_note={concepts: ideal gas equation, molar mass | steps: moles from PV equals nRT; molar mass of X2; atomic mass; digit sum}]
40 mg of gaseous substance (X$_2$) occupies 4.8 mL of volume at 1 atm and $27^\circ$C. Atomic mass of element X is : ($R=0.08$ atm L/mole-K)

(Fill the sum of digits of your answer in OMR sheet. e.g. if your answer is 1233 then fill the same as $1+2+3+3=9$)
\end{bankq}
\begin{bankq}[id=TOPS03-P2-CHEM-S1-Q01, source={TOPS - Test 03 (20 Nov 2022)}, paper=2, qno=1, type=multi, subtopic={Gas laws}, answer={A,D}, answer_src=key, difficulty=7, diff_parts={K2 S2 I2 C1 F1}, diff_note={concepts: ideal gas, graph interpretation | steps: ideal gas equation with the graph variables; slope and intercept; temperature values; intercept of graph 2 | log plot}]
Graph between $\log(1/n)$ v/s $\log P$ is plotted for an ideal gas enclosed (as shown in diagram) 24.63 litre container at three different temperatures.
\bfig{108.0}{TOPS03_p2_chem_s1_q01_1.png}
If $T_1=\dfrac{T_2}4=3T_3$ (where $T_1,T_2,T_3$ are temperature in kelvin of graph 1,2 and 3) then

($R=0.0821$ L-atm/K-mol $\log2=0.301$)
\begin{qoptions}
\opt{A}{$T_1$, $T_2$ and $T_3$ are 300, 1200 and 100 K respectively.}
\opt{B}{Slope of graphs 1, 2 and 3 are different.}
\opt{C}{Intercept of graph 1, 2 and 3 is same.}
\opt{D}{Intercept of graph 2 is 0.602.}
\end{qoptions}
\end{bankq}
\begin{bankq}[id=TOPS04-P1-CHEM-S1-Q01, source={TOPS - Test 04 (18 Dec 2022)}, paper=1, qno=1, type=multi, subtopic={Real gases}, answer={A,C}, answer_src=key, difficulty=5, diff_parts={K2 S2 I0 C1 F1}, diff_note={concepts: critical constants, van der Waals | steps: critical constants relations; Boyle temperature comparison; check each}]
Which of the following statements on critical constants of gasses are correct?
\begin{qoptions}
\opt{A}{Larger the value of $\dfrac{T_c}{P_c}$ of a gas, larger would be the excluded volume.}
\opt{B}{Critical temperature ($T_c$) of gas is greater than its Boyle temperature ($T_b$).}
\opt{C}{At the critical point in the van der Waals gas isotherm $\left(\dfrac{\partial p}{\partial V_m}\right)_{T_c}=0$}
\opt{D}{All are correct}
\end{qoptions}
\end{bankq}
\begin{bankpassage}[id=TOPS04-P1-CHEM-S3-G1-2, source={TOPS - Test 04 (18 Dec 2022)}, questions={1-2}]
Gas `A' is taken in a closed container at the initial total pressure of 1000 torr. Pressure of gas decreases to 900 torr in 5 seconds due to the effusion through a square cross-section. Another similar sized $2^{\text{nd}}$ container is taken in which gaseous mixture of $A$ and $B$ is taken at 5000 torr. Initial molar mass of mixture is $\dfrac{472}5$ gm/mol. A rectangular cross section is made in this container and gases are allowed to effuse. Width of this cross-section is same as the side of the previous square cross-section and length of the rectangular cross-section is 50\% more than that of its width.

(i) Gases $A$ and $B$ are non-reacting

(ii) Rate of effusion is directly proportional to area of orifice.

(iii) Rate of effusion is independent of the change in pressure due to diffusion.

Assume all other conditions to be identical. Answer the following question.

(Molar mass of $A=128$ gm/mole, Molar mass of $B=72$ gm/mole)
\end{bankpassage}
\begin{bankq}[id=TOPS04-P1-CHEM-S3-Q01, source={TOPS - Test 04 (18 Dec 2022)}, paper=1, qno=1, type=numerical, subtopic={Kinetic theory}, answer={33.33}, answer_src=key, difficulty=8, diff_parts={K2 S2 I2 C2 F1}, diff_note={concepts: effusion, mole fraction | steps: rate relation from the first container; rate ratio of the gases; composition from the molar mass; effused composition}, group=TOPS04-P1-CHEM-S3-G1-2]
mol \% of gas $A$ in mixture effused out initially from $2^{\text{nd}}$ container is :
\end{bankq}
\begin{bankq}[id=TOPS04-P1-CHEM-S3-Q02, source={TOPS - Test 04 (18 Dec 2022)}, paper=1, qno=2, type=numerical, subtopic={Kinetic theory}, answer={600.00}, answer_src=key, difficulty=8, diff_parts={K2 S2 I2 C2 F1}, diff_note={concepts: effusion, partial pressure | steps: rate of effusion A; size of the orifice; moles lost; pressure drop of B}, group=TOPS04-P1-CHEM-S3-G1-2]
Change is pressure (in torr) of gas $B$ in container in first 5 seconds due to effusion is-
\end{bankq}
\begin{bankq}[id=TOPS04-P1-CHEM-S4-Q07, source={TOPS - Test 04 (18 Dec 2022)}, paper=1, qno=7, type=numerical, subtopic={Real gases}, answer={4}, answer_src=key, difficulty=4, diff_parts={K1 S1 I1 C1 F1}, diff_note={concepts: ideal and real gas ideas | steps: check each of six statements; count the incorrect}]
Read the following statements :

(1) If a gas expands at constant temperature kinetic energy of the molecule is decreases.

(2) If a gas expands at constant pressure by providing heat, the kinetic energy of gaseous molecules remains same.

(3) At very low pressure real gases show minimum deviation from ideal behaviour.

(4) The compressibility factor for an ideal gas is zero.

(5) At boyle temperature real gas behave as ideal gas in high pressure region.

(6) Real gas show maximum deviation at high pressure and low temperature.

The number of incorrect statements is :
\end{bankq}
\begin{bankq}[id=TOPS04-P2-CHEM-S1-Q01, source={TOPS - Test 04 (18 Dec 2022)}, paper=2, qno=1, type=single, subtopic={Gas laws}, answer={A}, answer_src=key, difficulty=7, diff_parts={K2 S3 I2 C1 F0}, diff_note={concepts: gas mixing, reaction | steps: moles in each container; reaction of ammonia and HCl; remaining gas; final volume; pressure}]
Three containers are containers are connected by narrow tube as shown.
\bfig{133.7}{TOPS04_p2_chem_s1_q01_1.png}
Valve I is opened for long period of time and then closed. After it valve II is opened which results in immediate reaction as following :

$\mathrm{NH_3(g)+HCl(g)\longrightarrow NH_4Cl(s)}$

After long time valve II is closed. Calculate the final pressure (in atm) in container `B'.
\begin{qoptions}
\opt{A}{6 atm}
\opt{B}{7 atm}
\opt{C}{8 atm}
\opt{D}{9 atm}
\end{qoptions}
\end{bankq}
\begin{bankq}[id=TOPS04-P2-CHEM-S1-Q05, source={TOPS - Test 04 (18 Dec 2022)}, paper=2, qno=5, type=single, subtopic={Real gases}, answer={B}, answer_src=key, difficulty=5, diff_parts={K2 S2 I1 C1 F0}, diff_note={concepts: vapour pressure curves, phase | steps: identify the liquids from the curve; volatility; boiling condition}]
\bfig{82.0}{TOPS04_p2_chem_s1_q05_1.png}
(ii) For liquid $B$.

(iii) For liquid $C$.

Out of liquid $A$, $B$ \& $C$, identify correct option :
\begin{qoptions}
\opt{A}{Liquid $A$, $B$ \& $C$ can be water, ethanol, diethyl ether respectively.}
\opt{B}{Liquid $A$, $B$ \& $C$ can be diethyl ether, ethanol \& water respectively.}
\opt{C}{At temperature $T_1$, liquid $C$ is most volatile.}
\opt{D}{If temperature $>T_2$ \& pressure is 1 atm, then vapour phase of liquid $A$ will not exist.}
\end{qoptions}
\end{bankq}
\begin{bankq}[id=TOPS04-P2-CHEM-S2-Q02, source={TOPS - Test 04 (18 Dec 2022)}, paper=2, qno=2, type=numerical, subtopic={Real gases}, answer={2}, answer_src=key, difficulty=8, diff_parts={K2 S3 I2 C1 F1}, diff_note={concepts: compressibility factor, van der Waals | steps: Z in terms of b; slope at high pressure; solve b; excluded volume relation; molecular volume}]
The graph of compressibility factor $(Z)$ vs. $P$ for one mole of a real gas is shown in following diagram. The graph is plotted at constant temperature 273 K. If the slope of graph at very high pressure $\left(\dfrac{dZ}{dP}\right)$ is $\left(\dfrac1{2.8}\right)$ atm$^{-1}$, then calculate volume of one mole of real gas molecules (in L / mol)

Given : $N_A=6\times10^{23}$ and $R=\dfrac{22.4}{273}$ L atm K$^{-1}$mol$^{-1}$
\bfig{148.0}{TOPS04_p2_chem_s2_q02_1.png}
\end{bankq}
\begin{bankq}[id=TOPS04-P2-CHEM-S2-Q03, source={TOPS - Test 04 (18 Dec 2022)}, paper=2, qno=3, type=numerical, subtopic={Gas laws}, answer={3}, answer_src=key, difficulty=6, diff_parts={K2 S2 I1 C1 F1}, diff_note={concepts: ideal gas, process graph | steps: graph reading; moles and volume; partial derivative; value}]
Two moles of an ideal gas undergoes the following process. Given that $\left(\dfrac{\partial P}{\partial T}\right)_V$ is $x\times10^{-y}$ atm/k, then calculate the value of $(x+y)$
\bfig{129.6}{TOPS04_p2_chem_s2_q03_1.png}
\end{bankq}
\begin{bankq}[id=TOPS05-P1-CHEM-S1-Q02, source={TOPS - Test 05 (14 Jan 2023)}, paper=1, qno=2, type=multi, subtopic={Gas laws}, answer={A,C}, answer_src=key, difficulty=5, diff_parts={K1 S2 I1 C1 F1}, diff_note={concepts: ideal gas derivatives | steps: partial derivatives for one mole; product; simplify}]
For 1 mole of an ideal gas $\left(\dfrac{dP}{dT}\right)_V\cdot\left(\dfrac{dV}{dT}\right)_P=$
\begin{qoptions}
\opt{A}{$\dfrac RT$}
\opt{B}{$\dfrac1T$}
\opt{C}{$\dfrac{R^2}{PV}$}
\opt{D}{$\dfrac{R^2}{P^2}$}
\end{qoptions}
\end{bankq}
\begin{bankq}[id=TOPS05-P1-CHEM-S2-Q07, source={TOPS - Test 05 (14 Jan 2023)}, paper=1, qno=7, type=single, subtopic={Real gases}, answer={B}, answer_src=key, difficulty=5, diff_parts={K2 S2 I1 C1 F0}, diff_note={concepts: compressibility factor, Boyle temperature | steps: Z regions around Boyle temperature; forces dominant; match}]
{\footnotesize\setlength{\tabcolsep}{2pt}
\begin{tabular}{|l|p{3.1cm}|l|p{3.0cm}|}
\hline
 & List-I & & List-II\\
\hline
(P) & At Boyle's temperature in low pressure zone & (1) & $Z>1$\\
(Q) & At Boyle's temperature in high pressure zone & (2) & $Z=1$\\
(R) & Below Boyle's temperature in low pressure zone & (3) & Attractive forces are dominant between gas molecules\\
(S) & Above Boyle's temperature in high pressure zone & (4) & Repulsive forces are dominant between gas molecules\\
 & & (5) & $Z<1$\\
\hline
\end{tabular}}

Which is correctly matched :
\begin{qoptions}
\opt{A}{P$\to$5;Q$\to$1;R$\to$5;S$\to$1}
\opt{B}{P$\to$2;Q$\to$1;R$\to$5;S$\to$1}
\opt{C}{P$\to$3;Q$\to$1;R$\to$5;S$\to$2}
\opt{D}{P$\to$2;Q$\to$1;R$\to$4;S$\to$4}
\end{qoptions}
\end{bankq}
\begin{bankq}[id=TOPS05-P1-CHEM-S2-Q08, source={TOPS - Test 05 (14 Jan 2023)}, paper=1, qno=8, type=single, subtopic={Real gases}, answer={C}, answer_src=key, difficulty=3, diff_parts={K1 S1 I1 C1 F0}, diff_note={concepts: compressibility factor | steps: same matching; find the incorrect pair}]
Which is incorrectly matched :

{\footnotesize\setlength{\tabcolsep}{2pt}
\begin{tabular}{|l|p{3.1cm}|l|p{3.0cm}|}
\hline
 & List-I & & List-II\\
\hline
(P) & At Boyle's temperature in low pressure zone & (1) & $Z>1$\\
(Q) & At Boyle's temperature in high pressure zone & (2) & $Z=1$\\
(R) & Below Boyle's temperature in low pressure zone & (3) & Attractive forces are dominant between gas molecules\\
(S) & Above Boyle's temperature in high pressure zone & (4) & Repulsive forces are dominant between gas molecules\\
 & & (5) & $Z<1$\\
\hline
\end{tabular}}

Which is incorrectly matched ?
\begin{qoptions}
\opt{A}{P$\to$2;Q$\to$1,4}
\opt{B}{R$\to$3,5;S$\to$1,4}
\opt{C}{P$\to$2;Q$\to$1,3}
\opt{D}{Q$\to$1,4;R$\to$3,5}
\end{qoptions}
\end{bankq}

\section*{TMAS papers}
\begin{bankq}[id=TMAS01-P1-CHEM-S1-Q13, source={TMAS - Test 01 (09 Jul 2023)}, paper=1, qno=13, type=single, subtopic={Kinetic theory}, answer={B}, answer_src=key, difficulty=3, diff_parts={K1 S1 I1 C1 F0}, diff_note={concepts: effusion | steps: moles of each gas; ratio of rates}]
Certain amount of oxygen gas was passed over heated carbon, causing 80\% of oxygen added to be converted into CO gas, as per reaction\\
$\mathrm{2C(s)+O_2(g)\to2CO(g)}$\\
If the final gaseous mixture is allowed to effuse then ratio of rate of effusion of CO to $\mathrm{O_2}$ will be
\begin{qoptions}
\opt{A}{$\sqrt{\dfrac87}$}
\opt{B}{$\sqrt{\dfrac{8^3}7}$}
\opt{C}{$\sqrt{\dfrac{4^2\times8}7}$}
\opt{D}{$\sqrt{\dfrac{2^2\times8}7}$}
\end{qoptions}
\end{bankq}
\begin{bankq}[id=TMAS01-P2-CHEM-S1-Q04, source={TMAS - Test 01 (09 Jul 2023)}, paper=2, qno=4, type=multi, subtopic={Gas laws}, answer={A,B,C,D}, answer_src=key, difficulty=4, diff_parts={K1 S1 I1 C1 F1}, diff_note={concepts: ideal gas graphs | steps: relation for each plot; check each}]
Which of the following plot(s) is/are correct for a sample of ideal gas in given conditions ?
\begin{qoptions}
\opt{A}{\bfigi{72.0}{TMAS01_p2_chem_s1_q04_A.png}}
\opt{B}{\bfigi{119.3}{TMAS01_p2_chem_s1_q04_B.png}}
\opt{C}{\bfigi{74.5}{TMAS01_p2_chem_s1_q04_C.png}}
\opt{D}{\bfigi{97.5}{TMAS01_p2_chem_s1_q04_D.png}}
\end{qoptions}
\end{bankq}
\begin{bankq}[id=TMAS02-P2-CHEM-S1-Q02, source={TMAS - Test 02 (06 Aug 2023)}, paper=2, qno=2, type=multi, subtopic={Real gases}, answer={A,C}, answer_src=solved, working={$Z=1$ needs $0.2=60/T$, so $T=300$ K for any pressure}, difficulty=4, diff_parts={K1 S1 I1 C1 F1}, diff_note={concepts: compressibility factor | steps: condition Z equals 1; check each point}]
At moderate pressures, the compressibility factor for a particular gas is given by $Z=1+0.2P-\dfrac{60P}{T}$. Where $P$ is in bar and $T$ is in Kelvin. In which of the following conditions the gas behaves ideally :
\begin{qoptions}
\opt{A}{$P=1$ bar, $T=300$ K}
\opt{B}{$P=100$ bar, $T=400$ K}
\opt{C}{$P=2$ bar, $T=300$ K}
\opt{D}{$P=5$ bar, $T=330$ K}
\end{qoptions}
\end{bankq}
\end{document}
